\documentclass{article}
\usepackage[utf8]{inputenc}
\usepackage[T1]{fontenc}
\usepackage[margin=2cm]{geometry}
\usepackage{hyperref}
\sloppy
\begin{document}
\section*{Historical source}
Encyclopaedia Britannica scan, source row 1102. Parse 3.5 extracted the text. Super 3.5 VL corrected it from an enlarged crop of each block. It can still contain recognition errors.
Page headers, footers, and page numbers are removed.

\subsection*{Untitled continuation (E0091)}
\textit{Continues an entry from an earlier page.}

a matter of the utmost consequence to the success of the enterprise. Comines discharged this important duty with his usual ability, and delayed, though he could not prevent, the declaration of the republic in savour of the enemies of France; upon which he returned to Florence, where the king then was. Comines does not seem to have been employed in any public capacity during the last three years of Charles VIII.; but when Louis XII. ascended the throne in 1498, he appeared at court to render homage to the new sovereign. This is the last circumstance in the life of Comines of which any notice has been preserved. He died at Argenton, August 16, 1509, at the age of sixty-four. It was after his return from Italy that he commenced his Mémoires, at the solicitation of Angelo Cattho, Archbishop of Vienne, who, like himself, had been in the service first of the Duke of Burgundy, and afterwards in that of Louis XI.; and it is to this work that he is alone indebted for the celebrity of his name. Indeed, among modern historians, no one has been estimated more highly than Comines, who, to the charm of a natural and flexible style, which receives the direct impress of his thoughts, and to the interest which belongs to the fresh and animated recital of an eye-witness, unites a profound knowledge of men and affairs, and displays cool observation in happy alliance with sound judgment and accurate discrimination. In estimating the qualities of conduct and the shades of character, Comines is perhaps unrivalled. Taking men as he found them, and regarding them merely as the instruments of Providence, he depicts them with a cool, rigorous, yet apparently unstudied severity; and so true do his delineations appear, that he seems to have looked into their very souls, and to have discovered the real principles or motives of which their actions were the necessary result or expression. We look in vain in his Mémoires for noble sentiments, chivalrous feelings, or lofty aspirations. He never rises above the level of his subject, nor sacrifices that calm dignity of tone which he sustains without effort, because it seems to have been natural to his character. He speaks little of himself; and the most important circumstances of his life are omitted, probably because the recital of them might have embarrassed him. The first edition of the Mémoires of Comines, published at Paris in 1525, contains only the reign of Louis XI. from 1464 to 1483, and is divided into six books; but two more, containing the history of Charles VIII. till the year 1498, appeared in 1528, Paris, folio. A new edition, however, revised and corrected, was published at Paris in 1552, folio, by the Sieur du Parc, and has often since been reprinted. Of modern editions, the most complete and valuable is that of Lenglet-Dufresnoy, London, 1747, 4 vols. 4to.

Comines, or Commines, a town of Belgium, province of West Flanders, on the left bank of the Lys, 10 miles S.E. of Ypres. Pop. (1850) 3443. On the opposite bank of the river, in the department of Nord, is the French town of this name, with about 3000 inhabitants. The two towns are connected by a drawbridge. The manufactures are ribbons, thread, cottons, \&c.

\subsection*{COMITATUS (E0092)}
COMITATUS, in law, a county. England was first divided into counties, the counties into hundreds, and these again into tythings, by Alfred the Great. Sometimes it is taken for a territory or jurisdiction of a particular place.

\subsection*{COMITIA (E0093)}
COMITIA, in Roman Antiquity, general assemblies of the people, convened by a magistrate for the purpose of putting any subject to their vote.

The proper comitia were of three sorts, curiata, centuriata, and tributa, with reference to the three grand divisions of the city and people into curiae, centuriae, and tribes: for by comitia calata were meant all the comitia in general, the term calata (from calare, i. e., vocare) being their common epithet, though it was at last restricted to assemblies for the creation of priests, and the regulation of last wills and testaments, to which the people acted as witnesses merely.

The comitia curiata owed its origin to the division which Romulus made of the people into thirty curiae, ten being contained in every tribe. The curiae answered in most respects to the parishes in our cities, being not only separated by proper bounds and limits, but distinguished too by the different places set apart in each for the celebration of religious service, which was performed by particular priests (called curiones), one to every curia.

Before the institution of the comitia centuriata, all the great concerns of the state were transacted in the assembly of the curiæ; as the election of kings and other chief officers, the making and abrogating of laws, and the judging of capital crimes. After the expulsion of the kings, when the commons had obtained the privilege of having tribunes and ædiles, they were for some time elected at these assemblies; but that ceremony being at length transferred to the comitia tributa, the curiæ were never convened to give their votes, excepting upon occasion of making some particular law, relating to adoptions, wills, and testaments, or the creation of officers for an expedition, or for electing some of the priests, as the flamenes, and the curio maximus, who were themselves chosen by every particular curia. The power of calling these assemblies belonged at first only to the kings; but on the establishment of the democracy, this privilege was allowed to most of the magistrates, and sometimes to the pontifices.

The persons who had the liberty of voting were such Roman citizens as belonged to the curia, or actually lived in the city, and conformed to the customs and rites of their proper curia; and all were excluded who dwelt without the bounds of the city, or retained the ceremonies of their own country, even although they had been honoured with the jus civitatis. (See CITIZENS.) The place where the curia met was the comitium, a part of the forum. No fixed time was appointed for holding any of the comitia, but they were convened as business required.

The people being met together, and confirmed by the report of good omens from the augurs, which was necessary in all the assemblies, the rogatio or business to be proposed to them was publicly read. After this, if none of the magistrates interposed, upon the order of him who presided in the comitia the people divided into their proper curiae, and consulted of the matter; and then the curiae on being called out, as it happened by lot, gave their votes man by man, in ancient times viva voce, and afterwards by tablets; the most votes in every curia going for the voice of the whole curia, and the most curiae for the general consent of the people.

In the time of Cicero, the comitia curiata were so much out of fashion that they were formed only by thirty lictors representing the thirty curiae; whence, in his second oration against Rullus, he calls them comitia adumbrata.

The comitia centuriata were instituted by Servius Tullius, who had obliged every one to give a true account of what he was worth, and divided the people, according to the returns, into six ranks or classes, which he subdivided into a hundred and ninety-three centuries. The first class, containing the equites and richest citizens, consisted of ninety-eight centuries. The second, including the tradesmen and mechanics, consisted of twenty-two centuries; the third of twenty; the fourth of twenty-two; the fifth of thirty; and the sixth, composed of the poorer sort, of only one century; but this though it had the same name with the rest, yet was seldom regarded, and scarcely allowed any power in public matters. Hence it is common with the Roman authors, when they speak of the classes, to reckon no more than five, the sixth not being worth their notice. This last classis or order was divided into two parts or orders, the proletorii and the capite censi. The former, as their name implies, were designed purely to stock the republic with men, since they could supply it with so little money;


\section*{Source record}
\begin{description}
\item[Scan ID] \texttt{britannica:309f765b1ecc83e0a7f7762e44b4a927a7555443:train:1102}
\item[Image file] \texttt{page-1102-918dd1676cc2.jpg}
\item[Image SHA-256] \texttt{918dd1676cc2e6299e46d7d5610d5c3ee90a2c5fe401b45e30691f8e313293bc}
\item[Dataset] NationalLibraryOfScotland/encyclopaedia\_britannica\_illustrated at revision \texttt{309f765b1ecc83e0a7f7762e44b4a927a7555443} (CC0-1.0)
\item[Extraction] model \texttt{nvidia/nemotron-3.5-parse}, raw reply \texttt{extract/raw/page-1102-918dd1676cc2.json}
\item[Raw reply SHA-256] \texttt{d7ad6ff9a24ab5fe0e84755944ff896a9ea697637f95d0079bdee9fbfcc09f68}
\item[Correction] model \texttt{nvidia/Nemotron-3.5-Super-VL}, block crops, raw replies \texttt{correct/raw/crop\_nothink/page-1102-918dd1676cc2-b*.json}; super\_vl: 11
\item[Running header] 182 COM COM
\end{description}
\end{document}
